% =====================================================================
% Weld Agent Scope Study --- TOSEM paper skeleton (results pending)
%
% Target: ACM TOSEM, Special Section/Issue on Human-AI Collaboration in
% Software Engineering.
%   https://dl.acm.org/journal/tosem/continuous-special-sections-on-human-se
%   Submission type on Manuscript Central: "Human AI Collaboration in
%   Software Engineering" (handled as fast-impact track).
%
% BUILD:  latexmk -pdf WELD_TOSEM_PAPER.tex
%
% STATUS OF THIS FILE
%   Sections 3-5 (design, measures, analysis plan) are drafted as near-final
%   prose -- they are pre-specified and do not depend on the data.
%   Sections 1-2 and 6-9 are structured skeletons: each carries a \skel box
%   saying what goes in it, and every number that must come from data is a
%   \NUM{} placeholder so they can be grepped before submission.
%
% BEFORE SUBMITTING
%   - grep for \todo and \NUM -- none may survive.
%   - Verify every .bib entry below (drafted from memory; DOIs/venues/years
%     must be checked against the ACM DL / DBLP).
%   - Confirm the CfP's exact deadline, page limit, and whether the special
%     section requires double-anonymous review (the `anonymous' class option
%     below); the CfP PDF was not machine-fetchable at drafting time.
% =====================================================================

\documentclass[acmsmall,review,screen]{acmart}
% Add `anonymous' to the option list if the special section is double-anonymous.

\usepackage{booktabs}
\usepackage{tabularx}
\usepackage{array}

\newcolumntype{L}{>{\raggedright\arraybackslash}X}

% --- drafting aids (delete the definitions and all uses before submission) ---
\newcommand{\todo}[1]{\textcolor{red}{\textbf{[TODO:} #1\textbf{]}}}
\newcommand{\NUM}[1]{\textcolor{blue}{[#1]}}          % number that comes from data
\newcommand{\code}[1]{\texttt{#1}}
\newcommand{\skel}[1]{%
  \par\medskip\noindent
  \fbox{\parbox{0.96\linewidth}{\footnotesize\itshape\textbf{Skeleton note.} #1}}%
  \par\medskip}
\newcommand{\figplaceholder}[1]{%
  \fbox{\parbox[c][3.4cm][c]{0.92\linewidth}{\centering\itshape #1}}}

\begin{document}

\title{How Much Should an Agent See? A Controlled Study of Agent Scope
       in an Integrated Developer Workspace}
% Alternative titles to choose from:
%   - Seeing More, Doing More: Agent Scope and Developer Effectiveness
%   - Scope, Reliance, and Trust: A Controlled Study of a Workspace-Embedded
%     AI Agent
%   - What the Agent Can See: Manipulating Workspace Scope in Human-AI
%     Software Work

\author{Andrew Sadler}
\affiliation{%
  \institution{Michigan Technological University}
  \department{Department of Computer Science}
  \city{Houghton}\state{Michigan}\country{USA}}
\email{asadler@mtu.edu}

\author{Jie Wu}
\affiliation{%
  \institution{Michigan Technological University}
  \department{Department of Computer Science}
  \city{Houghton}\state{Michigan}\country{USA}}
\email{jie.jw.wu@mtu.edu}

\renewcommand{\shortauthors}{Sadler and Wu}

\begin{abstract}
\todo{Rewrite once results land. Structure below; keep to 200--250 words.}

AI assistants are increasingly embedded in the tools developers already work
in, yet how much of that workspace an assistant should be able to see and act
on remains an open design question. We report a controlled study
($N=\NUM{N}$) that manipulates exactly this variable. Participants worked in
Weld, a workspace that unifies a task board, team chat, and GitHub issues and
pull-request review, on a seeded software project. Each participant first
completed six tasks manually, establishing a per-person baseline, then completed
six matched tasks with an in-app AI agent whose scope over the workspace was
randomly assigned to one of three levels: chat only (advise), chat and board
(act), or chat, board, and GitHub (act). Task content was identical across
conditions; only the agent's scope differed. We logged task time, correctness,
delegation and override behaviour, and a per-invocation snapshot of exactly
what context the agent received, and collected NASA-TLX before and after the
agent phase together with a trust-in-automation block.
\NUM{Headline RQ1 result: where wider scope did and did not shorten
baseline-adjusted task time.}
\NUM{Headline RQ2 result: how delegation and verification shifted with scope.}
\NUM{Headline RQ3 result: whether trust tracked measured effectiveness.}
Our findings inform how much workspace access to grant assistants embedded in
software engineering tools, and \NUM{implication}.
\end{abstract}

\begin{CCSXML}
<ccs2012>
<concept>
<concept_id>10011007.10011074.10011099</concept_id>
<concept_desc>Software and its engineering~Software verification and validation</concept_desc>
<concept_significance>300</concept_significance>
</concept>
<concept>
<concept_id>10003120.10003121</concept_id>
<concept_desc>Human-centered computing~Human computer interaction (HCI)</concept_desc>
<concept_significance>500</concept_significance>
</concept>
<concept>
<concept_id>10003120.10003121.10003124.10010870</concept_id>
<concept_desc>Human-centered computing~Empirical studies in HCI</concept_desc>
<concept_significance>500</concept_significance>
</concept>
</ccs2012>
\end{CCSXML}

\ccsdesc[500]{Human-centered computing~Empirical studies in HCI}
\ccsdesc[300]{Software and its engineering~Software verification and validation}

\keywords{AI assistants, human-AI collaboration, controlled experiment,
developer tools, agent context, trust in automation, reliance}

\maketitle

% =====================================================================
\section{Introduction}
% =====================================================================

\skel{Five paragraphs, roughly one page. (1)~Software work is spread across
issue trackers, chat, and code hosts, and assistants are being embedded into
each of them. (2)~The open design question: assistants are increasingly given
wider access to that surrounding context, but the field has little controlled
evidence on what that access buys, or where. Prior work shows repository-level
context improves model output on benchmarks, and that developers' perceptions
of assistance diverge from measured behaviour --- but those are separate
literatures, tested separately, and neither manipulates a live assistant's
workspace scope with real users. (3)~This study: one workspace, one task set,
one variable --- how much of the workspace the agent can see and act on --- and
a per-person manual baseline. (4)~What we found. (5)~Contributions.}

\todo{Write after the results are in, so the framing matches the findings
rather than the hypotheses.}

\paragraph{Contributions.}
\begin{itemize}
  \item A controlled, pre-registered experiment isolating \emph{agent scope}
        over a shared developer workspace as an independent variable, with a
        within-person manual baseline and a per-invocation record proving the
        manipulation held.
  \item Evidence on \emph{where} scope matters: \NUM{result} across
        comprehension, single-surface, and cross-surface task types.
  \item Evidence on reliance: how delegation, verification, and override
        behaviour shift as the agent's scope widens.
  \item A characterisation of the perception--behaviour gap --- whether trust
        tracks measured effectiveness or the visible volume of agent activity.
  \item An open study apparatus (workspace, seeded project, task scripts,
        telemetry schema, scoring rubric) and anonymised data, enabling
        replication at other scope levels or with other models.
\end{itemize}

% =====================================================================
\section{Background and Related Work}
% =====================================================================

\skel{Four subsections, about 1.5 pages total. Each ends with the specific gap
this study fills --- not a survey. Cite the TOSEM SI's own framing where it
fits (human-AI collaboration, division of labour, trust and reliance).}

\subsection{AI assistants in software engineering practice}
\skel{Copilot-era usability and interaction studies: what developers do with
assistants, acceptance/rejection, validation and repair
\cite{vaithilingam2022expectation,barke2023grounded,liang2024largescale,mozannar2024reading}.
Gap: these study code completion inside an editor; almost none study an agent
acting across the \emph{project management} surfaces where much developer
coordination actually happens.}

\subsection{Context and agent capability}
\skel{Repository-level context improves generation and issue resolution
\cite{zhang2023repocoder,ding2023crosscodeeval,jimenez2024swebench}. Gap: these
are benchmark results about model output. Whether more context makes the
\emph{human} faster, more accurate, or more trusting is a separate question and
is not answered by benchmark gains.}

\subsection{Fragmentation, context switching, and coordination cost}
\skel{Developer time is fragmented across tools and interruptions carry a
measurable resumption cost \cite{meyer2014perceptions,meyer2017worklife,mark2008cost}.
Motivates the integrated workspace as apparatus. Be careful: this motivates the
setting, it does not license a claim that integration beats separate tools ---
we have no fragmented-tools arm (\S\ref{sec:threats}).}

\subsection{Trust, reliance, and the perception--behaviour gap}
\skel{Appropriate reliance and its failure modes
\cite{lee2004trust,parasuraman1997humans}; AI-assisted decision making where
explanations or visible activity raise trust without raising accuracy
\cite{bansal2021whole,bucinca2021trust}; measured-vs-perceived productivity for
coding assistants \cite{ziegler2022productivity}. Gap, and our RQ3: no prior
work varies an assistant's \emph{scope} and asks whether trust follows scope or
follows effectiveness. This is the study's most transferable finding and should
be foregrounded here.}

% =====================================================================
\section{Study Design}\label{sec:design}
% =====================================================================

\subsection{Research questions and hypotheses}

The study manipulates one variable: how much of the workspace an in-app AI
agent can see and act on. All questions stay within what that manipulation
supports.

\begin{description}
  \item[RQ1 (primary)] \textbf{Agent scope $\rightarrow$ effectiveness.} Does
    widening the agent's scope over the workspace (chat $\rightarrow$ +board
    $\rightarrow$ +GitHub) make the user more effective, and on which tasks?
  \item[RQ2] \textbf{Reliance.} How much do users delegate to versus verify
    and override the agent, and does that shift with its scope?
  \item[RQ3] \textbf{Trust versus reality.} Does users' trust in the agent
    track its measured effectiveness, or diverge from it?
\end{description}

\noindent Hypotheses:

\begin{description}
  \item[H1 (RQ1)] Wider scope yields shorter baseline-adjusted agent-task times
    \emph{across conditions}. Deliberately not stated as an absolute gain over
    the manual baseline: phase order is fixed, so an agent$-$baseline delta also
    contains ordinary practice (\S\ref{sec:timing}).
  \item[H2 (RQ1)] The gain is \emph{task-dependent} --- minimal on comprehension
    and single-surface tasks, substantial on cross-surface tasks spanning chat
    $\rightarrow$ board $\rightarrow$ GitHub.
  \item[H3 (RQ1)] Wider scope reduces self-reported cognitive load relative to
    baseline more than narrow scope does.
  \item[H4 (RQ2)] Delegation rises with scope, but reliance is imperfectly
    calibrated: users do not verify or override precisely where the agent is
    actually wrong.
  \item[H5 (RQ3)] Trust and perceived speedup need not track measured
    effectiveness --- users may trust an agent that \emph{does more} even when
    it is not faster, or is occasionally wrong.
\end{description}

\subsection{Design}

The design is hybrid: a within-subjects manual baseline followed by a
between-subjects agent condition.

Every participant first completes six manual baseline tasks (B1--B6, condition
C0) with no agent available, establishing their own manual speed and cognitive
load on each task type. NASA-TLX is collected at this point, before any agent
exposure. Each participant is then randomly assigned to one agent condition and
completes six agent tasks (T1--T6) in that fixed condition. Table~\ref{tab:conditions}
gives the three conditions.

\begin{table}[t]
\caption{The three agent scope conditions (between-subjects).}
\label{tab:conditions}
\begin{tabularx}{\linewidth}{@{}lLL@{}}
\toprule
\textbf{Condition} & \textbf{The agent can SEE} & \textbf{The agent can DO} \\
\midrule
C1 --- Chat only     & Chat                    & Advise only; the user executes \\
C2 --- Chat + board  & Chat, board             & Advise and act on board and chat \\
C3 --- Full scope    & Chat, board, GitHub     & Act across chat, board, and GitHub \\
\bottomrule
\end{tabularx}
\end{table}

\paragraph{Why hybrid.} The per-person C0 baseline lets each participant's
agent-task performance be analysed as a delta from their own manual speed,
absorbing individual skill differences that would otherwise sit between the
between-subjects conditions and sap power. Six baseline tasks --- one per agent
task type, four of them cross-surface --- give a type-matched, stable per-person
anchor rather than a single generic warm-up.

\paragraph{Why between-subjects on condition.} Holding the agent's access fixed
per person avoids within-person contamination: once a participant has seen a
full-scope agent solve a task, they cannot un-see how it should go. It also
keeps the session inside a self-serve 55--65 minute budget.

\paragraph{Capability scales with scope, by design.} An agent cannot act on a
board it cannot see, so visibility and capability are intentionally coupled.
This bundling is the phenomenon under study --- whether giving the agent more
access to the shared workspace helps the user --- and we therefore name the
variable \emph{agent scope} rather than \emph{context breadth} throughout. The
interpretive consequence is stated in \S\ref{sec:threats}.

\subsection{Apparatus}

\skel{One paragraph plus a screenshot figure. Describe Weld only as far as the
study needs: the surfaces, one seeded project, the agent panel. Do not turn this
into a tool paper.}

All participants work in \emph{Weld}, a workspace for coordinating and reviewing
software work: a Kanban task board, channel-based team chat, and GitHub
integration in a single application. Weld deliberately covers the coordination
and review layer around code rather than authoring or execution --- there is no
editor and no build --- which bounds what the study generalises to
(\S\ref{sec:threats}). GitHub is not a separate destination in the interface:
issues and pull requests are attached to the board cards they belong to, so a
participant reaches a pull request by opening its card, and reviews it in an
in-app review surface (description, files-changed diffs, and approve, comment,
or request-changes, plus merge). Merging moves the linked card exactly as a
GitHub webhook would in production use.

Every participant receives a private clone of the same fictional seed project,
``PetTracker'': 12 cards across 3 columns, 2 chat channels with 18 seeded
messages, and 3 pull requests, all 3 of which carry diffs. GitHub is mocked
from fixtures inside the application --- no live API, no
real repository, and outbound write controls hidden --- so participation is
fully self-serve and no participant's actions can affect another's environment
or leak into a real repository. Review verdicts and comments therefore exist
only in our event log, which is what makes them scorable
(\S\ref{sec:measures}).

The agent is presented as a panel inside the workspace. It runs on a single
pinned, dated model snapshot (\code{gpt-5.4-2026-03-05} --- matches the
\code{OPENAI\_MODEL} default in \code{config.rs}; confirm the deploy does not
override it via env var before the run starts) at a fixed
temperature of 0.2 with versioned prompts, held constant for the entire study.
A context gate assembles the agent's prompt per invocation under a default-deny
policy keyed to the participant's condition, and the same gate governs which
tools are exposed --- C1 has no tools at all, C2 may write to the board and post
to chat, and C3 may additionally link cards to issues and pull requests. Every
invocation records a snapshot of exactly what the agent received
(\S\ref{sec:manipcheck}).

\paragraph{The workspace re-presents GitHub; it does not replace it.} Issues,
pull requests, and diffs in Weld are a view onto data that lives in GitHub, and
we make no claim that this view is novel. That is deliberate on both sides. As
apparatus, a familiar GitHub-shaped review surface keeps the manipulation from
being confounded with learning an unfamiliar review tool. As a claim, what the
study varies is \emph{what the agent can reach}, not \emph{how it reaches it}:
the same contrast could be built by handing a C3 agent a GitHub API token or a
tool server, and we would expect the same result. The finding is therefore about
reach across the surfaces developer work spans, not about this particular
integration --- and it transfers to any assistant that can be granted or denied
that reach, however it is plumbed.

\begin{figure}[t]
\figplaceholder{Figure 1. The Weld workspace with the agent panel open.
Annotate the three surfaces (board, chat, GitHub) and show one agent turn.}
\caption{The study apparatus. \todo{Replace with a real screenshot; blur or
regenerate any content that could identify a participant.}}
\label{fig:apparatus}
\end{figure}

\subsection{Task set}

Twelve tasks --- six manual baselines then six agent tasks --- of roughly 4--5
minutes each. Task content is held constant across conditions; only the agent's
scope differs.

The agent tasks span three use-case types so we can locate \emph{where} scope
matters. Comprehension and single-surface tasks act as controls: they
demonstrate that wider scope does not \emph{universally} help, which is what
makes the cross-surface result credible rather than a blanket ``more is
better.'' Table~\ref{tab:tasks} summarises them; full scripts are in the
replication package.

\begin{table}[t]
\caption{The six agent tasks and what each is designed to test.}
\label{tab:tasks}
\begin{tabularx}{\linewidth}{@{}llL@{}}
\toprule
\textbf{Task} & \textbf{Type} & \textbf{What it tests} \\
\midrule
T1 ``What needs doing?''      & Comprehension  & Chat-only should be competitive \\
T2 ``Catch me up''            & Comprehension  & Do richer summaries reduce follow-up work? \\
T3 Draft a status message     & Single-surface (chat)  & Control: extra scope should not help \\
T4 Set priorities on the board& Single-surface (board) & Advising (C1) vs.\ acting (C2/C3) \\
T5 Bug from chat to the board & Cross-surface  & C2/C3 can act, C1 cannot --- core H1/H2 \\
T6 Review and link the PR, tell the team & Cross-surface & Only C3 has the scope to do it end to end \\
\bottomrule
\end{tabularx}
\end{table}

Each baseline task matches an agent task type --- B1 single-surface, B6
comprehension, B2--B5 cross-surface --- while using disjoint content, so it
anchors normalisation without teaching the agent tasks' answers.

T6 is constructed so the correct pull request is identifiable \emph{only} with
GitHub scope: chat refers to it vaguely, and the PR's title and body
disambiguate it. It is therefore the single cleanest test of the top of the
ladder --- if C3 does not beat C2 on T6, the GitHub scope is not pulling its
weight, and the data will say so. T6 also carries the most work: the
participant must identify the pull request, read its diff and record a review
verdict, connect it to the right card, and tell the team in chat. Its matched
baseline B4 requires the same review step on a different, unrelated pull
request, so the review cost is anchored per person rather than confounded with
the agent condition.

Task wording is load-bearing and was fixed in advance. Agent tasks state a goal
and never a route, and state that the assistant \emph{is available} rather than
instructing its use: naming the surfaces would let participants navigate without
the agent, collapsing the manipulation into UI mechanics, and instructing use
would turn RQ2's delegation measure into a compliance measure. T6 reads, in
full: \emph{``The vaccination reminders work isn't connected to its code yet,
and it hasn't been reviewed. Sort that out and let the team know in
\#pettracker. The assistant is available if you want it.''} Baseline tasks are
by contrast deliberately explicit about the route --- they are a per-person
normaliser, equally explicit for every condition, and their purpose is a stable
anchor rather than a test of navigation.

\subsection{Participants and procedure}

\skel{Fill the realised numbers from \code{study\_participants} plus the
demographics block; keep the planned targets in the pre-registration and report
deviations honestly.}

We target 20--30 participants (roughly 7--10 per condition), each contributing
six baseline and six agent observations. Eligibility: adults 18 or over, at
least one year of programming experience, and basic familiarity with issue
trackers, team chat, and GitHub. Recruitment is by self-serve link through
course lists, departmental channels, and developer communities; recruitment is
student-focused but not student-limited.

Sessions are unproctored and self-serve. The flow is: consent and eligibility
screening $\rightarrow$ a brief orientation to the workspace $\rightarrow$ six
manual baseline tasks $\rightarrow$ a post-baseline NASA-TLX $\rightarrow$ six
agent tasks in the assigned condition $\rightarrow$ the exit survey. Condition
assignment uses adaptive minimisation at provisioning time (each participant
goes to the least-filled condition, ties broken at random) rather than simple
randomisation, which keeps the cells balanced at this sample size.

Realised sample: $N=\NUM{N}$ (\NUM{n} C1, \NUM{n} C2, \NUM{n} C3),
\NUM{n} excluded (\NUM{reasons}). \todo{CONSORT-style participant flow diagram
from \code{study\_participants.status} and \code{session\_abandoned}.}

\subsection{Ethics}

\skel{One paragraph. State the IRB determination (number or exemption), consent
mechanics, the data-handling regime, and --- importantly --- the exclusion list,
which is a design commitment and not just boilerplate.}

This study was determined exempt from further review by Michigan Technological
University's Institutional Review Board. Participants consent before provisioning; consent is
recorded with the participant record. We deliberately do not collect keystroke
logs (only final submitted text), mouse movement or hover traces, screen or
camera recordings, data from other browser tabs, or any data about projects
other than the participant's own seeded instance. Data are pseudonymous
throughout and held in encrypted institutional storage; optional follow-up
contact details are stored separately, linked only by token, and deleted after
the follow-up window; open-ended text is screened for identifying content before
release.

% =====================================================================
\section{Measures}\label{sec:measures}
% =====================================================================

Every measure below names its source field, so the analysis is reproducible from
the event log and survey export. Table~\ref{tab:measures} summarises.

\begin{table}[t]
\caption{Measures and their sources.}
\label{tab:measures}
\begin{tabularx}{\linewidth}{@{}lLL@{}}
\toprule
\textbf{Construct} & \textbf{Source} & \textbf{Notes} \\
\midrule
Task completion time (primary) & \code{task\_completed.task\_time\_ms} &
  Active engaged time, not wall clock (\S\ref{sec:timing}). The clock stops
  before the per-task rating. \\
Away time, interruptions & \code{wall\_ms}, \code{away\_ms},
  \code{pause\_count}, \code{was\_interrupted} & Audit and exclusion flags. \\
Time to first action & first surface event after \code{task\_shown} & Derived. \\
Correctness (0--3) & manual coding of final state, agent transcript, DB &
  Two coders, blind to condition. \\
PR review behaviour & \code{pr\_review\_opened}, \code{pr\_review\_submitted}
  (verdict and comment), \code{pr\_merged} & B4/T6 only. Review comments are
  never written to GitHub, so this event is their only record and the rubric
  depends on it. \\
Per-task effort & \code{task\_rating.effort} (1--7, every task) &
  Within-person, baseline-anchored. \\
Global cognitive load & NASA-TLX, collected twice (post-baseline and exit) &
  The agent$-$baseline delta is the H3 measure. \\
Usability & SUS, \code{sus\_01}--\code{sus\_10} & Comparable to published norms. \\
Trust and reliance & 8-item trust block & H5; \code{trust\_verify} is
  reverse-valenced. \\
Agent effectiveness & \code{agent\_response} (latency, response type, tool
  calls, tokens), \code{agent\_tool\_call} & Effort and cost. \\
Delegation & executed agent tool calls vs.\ user actions; accept / reject /
  edit / undo events; \code{agent\_panel\_opened} & RQ2. Panel-open is the
  coarsest measure of all: did the participant reach for the agent? \\
Cross-reference usage & card-from-chat, PR-from-card, PR-from-chat,
  card-from-PR, issue-from-card & Integration-usage measure. \\
View-switch rate & \code{view\_changed} count per unit active time &
  Context-switching between board and chat only --- see below. \\
Scoping proof & \code{agent\_context\_sent} & Manipulation check
  (\S\ref{sec:manipcheck}). \\
\bottomrule
\end{tabularx}
\end{table}

\paragraph{Baseline-normalised time.} The headline derived measure is each agent
task's active time relative to that participant's own type-matched baseline:
comprehension tasks (T1, T2) against B6, single-surface tasks (T3, T4) against
B1, and cross-surface tasks (T5, T6) against the mean of B2--B5. Every agent
task type therefore has a like-for-like anchor, and the cross-surface anchor ---
where the hypotheses live --- averages four baselines and is the most stable. In
the models the baseline enters as a covariate rather than as a pre-computed
ratio (\S\ref{sec:models}).

\paragraph{What the view-switch rate does and does not cover.} Because GitHub is
reached from the card it belongs to rather than from a separate destination, a
route-level view change fires for board, chat, and settings navigation but never
for opening a pull request. The view-switch rate is therefore a board--chat
context-switching measure, and traversal into GitHub is captured by the
cross-reference events and \code{github\_view\_opened} instead. The two are
reported together rather than summed, and the perceived-switching survey item is
paired with the board--chat rate only. \todo{Confirm this is still true of the
build used in the run --- it is a consequence of a UI decision, not a study
setting, so a later interface change silently changes the measure.}

\paragraph{Correctness.} Speed without correctness is not effectiveness, so
correctness is scored for every task on a 0--3 rubric and always reported beside
time. Scoring procedure is fixed in advance because it is the one primary
measure produced by judgement rather than instrumentation: two coders score all
tasks from a final-state export with condition, agent transcript, and timing
stripped --- blinding is mandatory, since knowing a participant was in C3 while
judging their output is a direct route to confirming H1. We report
quadratic-weighted Cohen's $\kappa$ \emph{before} adjudication, then resolve
disagreements to consensus, still blind. The rubric is piloted and frozen before
any real session is scored.

% =====================================================================
\section{Analysis Plan}\label{sec:analysis}
% =====================================================================

This plan was fixed before data collection. Deviations are reported explicitly
in \S\ref{sec:results}.

\subsection{Manipulation check}\label{sec:manipcheck}

Reported first, before any inferential result. From the \code{agent\_context\_sent}
snapshot recorded at every invocation, we verify that 100\% of C1 snapshots
contain no board or GitHub content and 100\% of C2 snapshots contain no GitHub
content. This converts the independent variable from an assumption into
evidence; without it the scoping manipulation rests on trust.

\subsection{Primary models}\label{sec:models}

RQ1/H1/H2 are tested with a linear mixed model on log task time,
\[
\log(\mathit{time}) \sim \mathit{condition} \times \mathit{task\_type}
  + \log(\mathit{baseline}) + (1 \mid \mathit{participant}),
\]
with H2 as the \code{condition}~$\times$~\code{task\_type} interaction: does the
condition effect concentrate on the cross-surface tasks? At this sample size the
full interaction is underpowered, so the confirmatory analysis is a
pre-registered focused contrast --- C3 versus C1 on cross-surface tasks --- with
the omnibus interaction reported as supporting evidence and everything else
flagged exploratory. We lead with effect sizes and 95\% confidence intervals
rather than $p$-values throughout.

Correctness is modelled in parallel with an ordinal cumulative-link mixed model,
\code{correctness \textasciitilde{} condition * task\_type + (1 | participant)},
and reported beside every timing result.

\subsection{Cognitive load (H3)}

NASA-TLX is collected twice, so load is a within-person phase-to-phase delta
rather than a single end-of-session snapshot. The headline test is
$\Delta\mathrm{TLX} = \mathrm{TLX}_{\text{agent}} - \mathrm{TLX}_{\text{baseline}}$
on the composite and per subscale, compared between conditions (Welch ANOVA or
Kruskal--Wallis given the sample size). A finer-grained secondary analysis models
per-task effort with each participant's baseline effort as the anchor.

\subsection{Reliance (RQ2) and the perception--behaviour gap (RQ3)}

RQ2 is answered descriptively from telemetry: delegation ratio by condition; the
agent funnel from invocation through attempted and executed tool calls to
accepted, edited, or undone outcomes; tool error and rejection rates; and
cross-reference traversal counts. H4's calibration claim is tested by crossing
override behaviour against coded agent-error type --- do users override where the
agent is actually wrong?

RQ3 correlates the trust composite and \code{trust\_speed\_perceived} (Spearman)
against measured normalised time and correctness, and characterises any
divergence. \code{trust\_rely} crossed with \code{trust\_verify} gives reliance
calibration; perceived context-switching is paired with the behavioural
view-switch rate. Four open-ended items are thematically coded by two coders with
$\kappa$ reported on a shared subset.

\subsection{Timing validity}\label{sec:timing}

Task time is \emph{active engaged time}: the client clock runs only while the tab
is visible, pauses on background or close, resumes on return, and persists across
reloads, so away-time is excluded by construction rather than after the fact.
Each task additionally logs away time, pause count, and an interruption flag. We
analyse active time as primary, run a sensitivity analysis excluding interrupted
tasks, and report the interruption rate.

Because the clock catches stepping away from the tab but not stepping away from a
tab left in front, one further rule was pre-specified: within a task, 180 seconds
or more of visible time containing zero telemetry events is treated as idle and
subtracted. The threshold sits well above a natural pause for reading a channel
or a diff --- tasks run 4--5 minutes and emit events throughout --- and well
below a genuine walk-away. We report how many tasks trigger it and fit the
primary model with and without the correction.

Phase order is fixed: the six manual baselines always precede the six agent
tasks, since agent exposure cannot be undone. An agent$-$baseline delta therefore
contains ordinary practice as well as agent help. Between-condition inferences
are unaffected --- every arm has the identical ordering --- but we never claim
the agent beat manual work.

% =====================================================================
\section{Results}\label{sec:results}
% =====================================================================

\skel{Results go here in this order --- manipulation check first, then RQ1, RQ2,
RQ3. Every subsection is currently a placeholder naming the exact table or
figure it will contain. Report the realised sample and any deviation from
\S\ref{sec:analysis} at the top.}

\subsection{Manipulation check}

\NUM{Table: per-condition counts of agent invocations and the proportion of
context snapshots containing board / GitHub content. Expect 0\% off-scope.}
\todo{If any snapshot leaked off-scope context, report it here and state how the
affected observations were handled --- do not quietly drop them.}

\subsection{RQ1 --- Agent scope and effectiveness}

\NUM{Figure 2: normalised time by condition $\times$ task type --- the H2
interaction plot.}
\NUM{Figure 3: correctness on the same axes.}
\NUM{Table: model coefficients with 95\% CIs, plus cell-level Cohen's $d$.}
\NUM{The pre-registered focused contrast, C3 vs.\ C1 on cross-surface tasks.}
\NUM{T6 specifically: C3 vs.\ C2, the cleanest top-of-ladder test.}

\subsection{RQ2 --- Reliance and the division of labour}

\NUM{Figure 4: delegation ratio by condition.}
\NUM{Table: the agent funnel --- invoked, tool attempted, executed, accepted,
edited, undone --- by condition.}
\NUM{Calibration: override rate crossed with coded agent-error type.}

\subsection{RQ3 --- Trust versus measured effectiveness}

\NUM{Figure 5: perception--behaviour scatter, trust composite against measured
normalised time, points coloured by condition.}
\NUM{Correlations with CIs; trust\_rely $\times$ trust\_verify cross-tab.}

\subsection{Cognitive load and usability}

\NUM{$\Delta$TLX composite and subscales by condition; SUS score against
published norms.}

\subsection{Qualitative findings}

\NUM{Themes from the four open-ended items with $\kappa$; two or three
representative quotes per theme, screened for identifying content.}

% =====================================================================
\section{Discussion}
% =====================================================================

\skel{Three or four subsections, each anchored to a result. Resist writing this
before the data --- the placeholders below are angles, not conclusions.}

\subsection{Where scope pays, and where it does not}
\skel{If the interaction holds: scope is not a monotone good, and tool builders
should widen an assistant's access along the seams the work actually crosses
rather than uniformly. If it does not hold: say so plainly, and use the T6
contrast to say whether the ceiling was the scope or the model.}

\subsection{Implications for designers of workspace-embedded assistants}
\skel{Concrete, actionable: what to grant, what to gate, what to surface to the
user about what the assistant can see. Tie back to the TOSEM SI's
human-AI-collaboration framing --- this is the section the special issue is
asking for.}

\subsection{Trust as a design risk}
\skel{If trust tracks visible activity rather than effectiveness, an assistant
that does more is trusted more regardless of whether it is right --- a design
hazard, since widening scope is exactly what makes an assistant visibly do
more. Connect to appropriate-reliance work \cite{lee2004trust,bucinca2021trust}.}

\subsection{Implications for research}
\skel{Benchmark context gains do not transfer automatically to user outcomes;
name what a future study with a fragmented-tools arm, or with capability held
fixed across scope levels, would need to look like.}

% =====================================================================
\section{Threats to Validity}\label{sec:threats}
% =====================================================================

\skel{Write these honestly and early --- several are known now and are stronger
stated up front than extracted by a reviewer.}

\paragraph{Scope confounds seeing with doing.} Visibility and capability rise
together by construction, so the C3-versus-C1 contrast confounds seeing with
doing: C1 only advises while C3 acts across three surfaces, and the simplest
reading of any gap is that the agent did the clicking. We report that contrast as
evidence the manipulation works, not as evidence about context per se. The
contrast that isolates context is C2 versus C3 --- both act, and differ only in
GitHub scope --- which this sample can estimate but not decisively test; T6
sharpens it. RQ3 is untouched by this confound, since it compares each
condition's trust against that same condition's own measured effectiveness.

\paragraph{No fragmented-tools arm.} The study cannot speak to whether an
integrated workspace beats separate tools. Integration motivates the apparatus;
it is not a claim we test.

\paragraph{The top rung is small, even though it is real.} Participants do
genuinely review in C3 --- they read a diff, record a verdict and a comment,
link the pull request, and can merge it, with the merge moving the linked card
as a webhook would. But the diffs are seeded fixtures of one to two files, so
review here costs a fraction of what review costs on a real change set. That
caps the detectable C2$\rightarrow$C3 effect --- if reviewing is cheap, help
with reviewing is worth little --- and it limits generalisation to real
repositories, where the reading is the expensive part and the effect could
plausibly be larger. \todo{After the run, report the actual median time spent in
the review surface; if it is very short, say so here rather than letting a
reviewer compute it.}

\paragraph{Coordination, not authoring.} Weld covers the coordination and review
layer around code. Participants never write or run code, so the findings speak to
agents embedded in the tools \emph{around} an editor, not to in-editor coding
assistants, which is where most existing evidence sits.

\paragraph{One model, one moment.} A single pinned snapshot at fixed temperature
buys internal validity at the cost of external validity. The snapshot must also
be capable enough that the top rung is floored by the design rather than by the
model; we verified this in piloting.

\paragraph{Fixed phase order.} Practice effects are bundled into any
agent$-$baseline delta (\S\ref{sec:timing}), which is why H1 is stated as a
between-condition difference.

\paragraph{Sample size and setting.} Roughly 7--10 participants per cell, on a
seeded fictional project, in unproctored self-serve sessions, with a
student-weighted sample. \todo{State the realised sample's demographics and be
specific about what it does and does not generalise to.}

% =====================================================================
\section{Conclusion}
% =====================================================================

\skel{Half a page. Restate the manipulation, the headline finding, and the one
design lesson worth carrying away. Write last.}

\todo{Write after the discussion.}

\section*{Data Availability}

The study apparatus (workspace source, seed project, task scripts, telemetry
schema, correctness rubric, and survey instrument) and the anonymised data and
analysis scripts are available at \todo{archival DOI}.

\bibliographystyle{ACM-Reference-Format}
\bibliography{weld-tosem}

\end{document}
